% ============================================================
% CHAPTER 1: INTRODUCTION
% ============================================================
\chapter{Introduction}
\label{ch:introduction}

Cycling has been widely recognised as a sustainable mode of urban transport, offering co-benefits across public health, air quality, carbon reduction, and potential congestion alleviation \parencite{who2022,ipcc2023,eu2024}. Central to realising these benefits is the provision of high-quality, connected cycling infrastructure, which is consistently associated with higher cycling participation \parencite{hull2014,pistoll2014,buehler2016,molenberg2019}. Reflecting this evidence, the UK Department for Transport published Local Transport Note 1/20, Cycle infrastructure design (LTN 1/20), in July 2020, establishing the most comprehensive national standards for cycling infrastructure design and assessment applicable to England and Northern Ireland to date \parencite{departmentfortransport2020a}. LTN 1/20's evaluation tools, the Cycling Level of Service (CLoS) and the Junction Assessment Tool (JAT), give practitioners a rigorous means of appraising individual schemes, assessing road segments and junctions against detailed design criteria through structured manual auditing \parencite{departmentfortransport2020a}. What they do not address is the network-level question: whether compliant segments cohere into a connected whole. Network-level cycling connectivity is a property of whole journeys rather than individual schemes: a single high-stress link or junction can render an otherwise compliant route unusable for the very cyclists the guidance seeks to serve. There is therefore a need for a data-driven, network-level methodology capable of evaluating cycling infrastructure quality across a city or region using standard open and mapping agency data.

The Level of Traffic Stress (LTS) framework offers such a network-level approach, classifying road segments and junctions by the population group willing to tolerate their traffic conditions: from LTS 1, conditions suitable for children and old people, through LTS 2, tolerable to most adults, to LTS 3 and 4, which only progressively smaller and more confident groups of experienced cyclists will accept \parencite{mekuria2012a}. Because each stress level maps onto a population group, the classification links infrastructure quality directly to mode shift potential: expanding the network usable at lower stress thresholds is what brings the broader population into cycling. LTS-based network analysis has since become a leading approach in cycling connectivity and equity research \parencite{louro2026}. Its non-compensatory logic, under which a single high-stress attribute determines classification, aligns naturally with how traffic conditions deter potential cyclists. Yet most published LTS criteria were calibrated to North American road design conventions, and applications to British cities have relied on simplified adaptations of these criteria rather than national design guidance \parencite{cervero2019,jeong2025}. No LTS classification grounded in UK design standards has been developed systematically.

The recent release of the Ordnance Survey National Geographic Database (OS NGD) Cycle Lane feature collection provides, for the first time, a nationally consistent and authoritative cycling infrastructure dataset, with facility descriptions aligned with LTN 1/20 terminology. The wider OS NGD Transport Theme complements this by providing network geometry together with attributes directly relevant to stress classification, including carriageway width, legal access designation and probe-derived vehicle speeds. The Cycle Lane feature collection, which achieved full national coverage in March 2026, additionally records cycle lane width, direction of cycling travel and infrastructure type \parencite{ordnancesurvey2026e}. Together, the availability of a statutory design standard and a nationally consistent dataset enables cycling networks across England to be classified and evaluated against a single national benchmark.

This dissertation develops a directed Level of Traffic Stress classification derived from LTN 1/20 design criteria, operationalises it on OS NGD Transport data, and applies it to evaluate network-level cycling conditions in the West Midlands. The region presents a revealing setting for such an evaluation. Among English regions it combines the highest share of private motorised trips with a cycling mode share of only 0.9\% \parencite{westmidlandscombinedauthority2024}, yet its transport authority has committed to delivering a 500-mile network of low-stress cycle routes \parencite{westmidlandscombinedauthority2020}. The distance between this ambition and current network conditions is precisely what a standards-based classification is designed to measure. The dissertation, conducted in collaboration with Transport for West Midlands, makes contributions to methodology and data: the first LTS framework calibrated to English design standards systematically, incorporating directed junction assessment and roundabout-specific criteria; and the first independent assessment of the OS NGD Cycle Lane dataset. The classified network is then used to evaluate connectivity against LTN 1/20's own principle of “cycling for all”, proceeding from infrastructure supply to the journeys that infrastructure enables: the extent and spatial distribution of low-stress provision, the degree to which low-stress links cohere into connected components rather than isolated fragments, the reach available to residents from where they live, and the proportion of everyday travel demand, for access to schools and for commuting, that the network can carry at each stress threshold. A network growth simulation then examines how the sequence in which high-stress links are upgraded shapes connectivity outcomes, comparing selection strategies that exploit network structure against uninformed investment. Together these analyses illustrate the range of strategic questions that a nationally consistent, standards-based classification can support.
